\Needspace{8\baselineskip}
\Article{208}
\textsc{Problem.} \textit{Given two properly equivalent forms $F$, $F'$ with determinant $hh$, find a proper transformation of one into the other.}

\SourceLeaf{207}
\textit{Solution.} Let $F$ be properly equivalent to a reduced form $\Phi$, which, by hypothesis, is therefore also properly equivalent to $F'$. By Article~206, find a proper transformation of $F$ into $\Phi$, say $\alpha,\beta,\gamma,\delta$, and also a proper transformation of $F'$ into $\Phi$, say $\alpha',\beta',\gamma',\delta'$. Then $\Phi$ passes into $F'$ by the proper substitution $\delta',-\beta',-\gamma',\alpha'$, and so $F$ passes into $F'$ by the proper substitution
\[
\alpha\delta'-\beta\gamma',\quad\beta\alpha'-\alpha\beta',\quad
\gamma\delta'-\delta\gamma',\quad\delta\alpha'-\gamma\beta'
\]

It is useful to develop another formula for this transformation of $F$ into $F'$, for which we need not even know the reduced form $\Phi$ itself. Let
\[
F\text{ be }(a,b,c),\quad F'=(a',b',c'),\quad\Phi=(A,h,0)
\]
Since the ratio in smallest terms equal to $h-b:a$ or $c:-(h+b)$ is $\beta:\delta$, we easily see that $\frac{h-b}{\beta}=\frac a\delta$ is an \textit{integer}; call it $f$. Likewise, $\frac c\beta=\frac{-h-b}{\delta}$ is an integer; put it $=g$. We have
\[
\begin{gathered}A=a\alpha\alpha+2b\alpha\gamma+c\gamma\gamma\quad\text{so}\\
\beta A=a\alpha\alpha\beta+2b\alpha\beta\gamma+c\beta\gamma\gamma\end{gathered}
\]
or, substituting for $a\beta$ the value $\delta(h-b)$, and for $c$ the value $\beta g$,
\[
\beta A=\alpha\alpha\delta h+b(2\beta\gamma-\alpha\delta)\alpha+\beta\beta\gamma\gamma g
\]
or, since $b=-h-\delta g$,
\[
\beta A=2\alpha(\alpha\delta-\beta\gamma)h+(\alpha\delta-\beta\gamma)^2g=2\alpha h+g
\]
Similarly,
\[
\begin{aligned}
\delta A&=a\alpha\alpha\delta+2b\alpha\gamma\delta+c\gamma\gamma\delta\\
&=\alpha\alpha\delta\delta f+b(2\alpha\delta-\beta\gamma)\gamma-\beta\gamma\gamma h\\
&=(\alpha\delta-\beta\gamma)^2f+2\gamma(\alpha\delta-\beta\gamma)h=2\gamma h+f
\end{aligned}
\]
So
\[
\alpha=\tfrac{\beta A-g}{2h},\qquad\gamma=\tfrac{\delta A-f}{2h}
\]
In exactly the same way, putting
\[
\frac{h-b'}{\beta'}=\frac{a'}{\delta'}=f',\qquad
\frac{c'}{\beta'}=\frac{-h-b'}{\delta'}=g'
\]
gives
\[
\alpha'=\tfrac{\beta' A-g'}{2h},\qquad\gamma'=\tfrac{\delta' A-f'}{2h}
\]
\SourceLeaf{208}
Substituting these values of $\alpha,\gamma,\alpha',\gamma'$ into the formula just given for the transformation of $F$ into $F'$ turns it into
\[
\tfrac{\beta f'-\delta'g}{2h},\quad
\tfrac{\beta'g-\beta g'}{2h},\quad
\tfrac{\delta f'-\delta'f}{2h},\quad
\tfrac{\beta'f-\delta g'}{2h}
\]
from which $A$ has disappeared entirely.

If two improperly equivalent forms $F$, $F'$ are given and we seek an improper transformation of one into the other, let $G$ be the form opposite to $F$, and let a proper transformation of $G$ into $F'$ be $\alpha,\beta,\gamma.\ \delta$. Then clearly $\alpha,\beta,-\gamma,-\delta$ is an improper transformation of $F$ into $F'$.

Finally, if the given forms are both properly and improperly equivalent, we can clearly find two transformations in this way, one proper and one improper.

\Needspace{8\baselineskip}
\Article{209}
All that remains is to show how to derive all the other similar transformations from one transformation. This depends on solving the indeterminate equation $tt-hhuu=mm$, where $m$ denotes the greatest common divisor of $a,2b,c$, if $(a,b,c)$ is either of the equivalent forms. But this equation can always be solved in only two ways: either put $t=m$, $u=0$, or $t=-m$, $u=0$. For suppose there is another solution $t=T$, $u=U$, with $U$ not $=0$. Since $mm$ certainly divides $4hh$, we have $\frac{4TT}{mm}=\frac{4hhUU}{mm}+4$, and both $\frac{4TT}{mm}$ and $\frac{4hhUU}{mm}$ are integer squares. But we easily see that $4$ cannot be the difference of two integer squares unless the smaller square is $0$, \textit{i.e.}, $U=0$, contrary to the hypothesis. --- So if $F$ passes into $F'$ by the substitution $\alpha,\beta,\gamma,\delta$, there is no other transformation similar to this one except $-\alpha,-\beta,-\gamma,-\delta$. If two forms are equivalent only properly, or only improperly, there are therefore only \textit{two} transformations. If they are both properly and improperly equivalent, there are \textit{four}, namely two proper and two improper.

\Needspace{8\baselineskip}
\Article{210}
\textsc{Theorem.} \textit{If two reduced forms $(a,h,0)$, $(a',h,0)$ are improperly equivalent, then $aa'\equiv mm\;(\mathit{mod.}\,2mh)$, where $m$ denotes the greatest common divisor of $a,2h$ or $a',2h$. Conversely, if $a,2h;\ a',2h$ have the same greatest common}
\SourceLeaf{209}
\textit{divisor $m$ and $aa'\equiv mm\;(\mathit{mod.}\,2mh)$, then the forms $(a,h,0)$, $(a',h,0)$ are improperly equivalent.}

\textit{Proof.} I. Let $(a,h,0)$ pass into $(a',h,0)$ by the improper substitution $\alpha,\beta,\gamma,\delta$, so that we have four equations:
\begin{align*}
a\alpha\alpha+2h\alpha\gamma&=a'\tag*{. . . . . . . . . . . . [1]}\\
a\alpha\beta+h(\alpha\delta+\beta\gamma)&=h\tag*{. . . . . . . . . . [2]}\\
a\beta\beta+2h\beta\delta&=0\tag*{. . . . . . . . . . . . [3]}\\
\alpha\delta-\beta\gamma&=-1\tag*{. . . . . . . . . . . . . [4]}
\end{align*}
Multiplying [4] by $h$ and subtracting it from [2], which we write as $[2]-h[4]$, gives
\begin{equation*}
(a\alpha+2h\gamma)\beta=2h\tag*{. . . . . . . . . . . . [5]}
\end{equation*}
Similarly, from $\gamma\delta[2]-\gamma\gamma[3]-(a+a\beta\gamma+h\gamma\delta)[4]$, after removing the terms that cancel, we get
\begin{equation*}
-a\alpha\delta=a+2h\gamma\delta,\quad\text{or}\quad -(a\alpha+2h\gamma)\delta=a\tag*{. . [6]}
\end{equation*}
Finally, we get $a[1]\ldots a\alpha(a\alpha+2h\gamma)=aa'$, or
\[
(a\alpha+2h\gamma)^2-aa'=2h\gamma(a\alpha+2h\gamma)
\]
or
\begin{equation*}
(a\alpha+2h\gamma)^2\equiv aa'\;(\mathrm{mod.}\,2h(a\alpha+2h\gamma))\tag*{. . . . . [7]}
\end{equation*}
Now [5] and [6] show that $a\alpha+2h\gamma$ divides $2h$ and $a$, and therefore also $m$, the greatest common divisor of $a,2h$. But clearly $m$ also divides $a\alpha+2h\gamma$. So $a\alpha+2h\gamma$ must be either $=+m$ or $=-m$. It follows at once from [7] that $mm\equiv aa'\;(\mathrm{mod.}\,2mh)$. \textit{This proves the claim.}

II. If $a,2h;\ a',2h$ have the same greatest common divisor $m$ and, in addition, $aa'\equiv mm\;(\mathrm{mod.}\,2mh)$, then $\frac a m$, $\frac{2h}{m}$, $\frac{a'}{m}$, $\frac{aa'-mm}{2mh}$ are integers. It is easy to check that $(a,h,0)$ passes into $(a',h,0)$ by the substitution $-\frac{a'}{m},-\frac{2h}{m},\frac{aa'-mm}{2mh},\frac a m$, and that this transformation is improper. So these forms are improperly equivalent. \textit{This proves the claim.}

From this we can also decide at once whether a given reduced form $(a,h,0)$ is improperly equivalent to itself. If $a,2h$ have greatest common divisor $m$, the required condition is $aa\equiv mm\;(\mathrm{mod.}\,2mh)$.

\SourceLeaf{210}
\Needspace{8\baselineskip}
\Article{211}
We obtain all reduced forms with given determinant $hh$ by substituting in the indefinite form $(A,h,0)$, for $A$, all numbers from $0$ to $2h-1$, inclusive. So there are $2h$ of them. Clearly all forms with determinant $hh$ can be divided into that many \textit{classes}, and these have the same properties we discussed above (Articles~175, 195) for classes of forms with negative determinant and positive nonsquare determinant. For example, all forms with determinant $25$ divide into ten classes, which can be distinguished by the reduced form in each. These reduced forms are $(0,5,0)$, $(1,5,0)$, $(2,5,0)$, $(5,5,0)$, $(8,5,0)$, $(9,5,0)$, which are also improperly equivalent to themselves; $(3,5,0)$, to which $(7,5,0)$ is improperly equivalent; and $(4,5,0)$, to which $(6,5,0)$ is improperly equivalent.

